% Chapter 6, section 5: Consumption engine
% Write in plain LaTeX. Code of a guide step: \guide{<step>}; commands: \begin{shell}...\end{shell};
% a screenshot: \screenshot{file}{caption}; a gap in the draft: \placeholder{...}.
% See book/latex/queen.sty for the full list.

\section{Consumption engine}

Pops, acks and leases do not go through the planner. The leader keeps each consumer group's
positions and leases in memory, and gives out the oldest ready partition first.

The positions that changed go into the log as one small checkpoint, at once when an answer is
waiting and at least every 5~ms. A pop gets its answer only after that checkpoint is committed. A
pop that finds nothing waits, and is woken when a new message arrives in its partition.

\queenfig{consumption-engine}{Pops, acks and leases, served from the leader's memory.}

We moved consumption out of the planner because acks were waiting behind pushes: at 300,000
messages a second, an ack took more than 500~ms. Now many acks travel in one checkpoint.

A partition idle for 10 minutes leaves memory, and is loaded again with its next message. A new
leader waits about half a second before it serves pops, so the same messages are never given out
twice.
